% GENERATED FILE. Edit canonical lessons, then run npm run generate:books.
% canonical-source-hash: fnv1a64:925a43adcddfb101
% canonical-lessons: TE-C310-parsu, TE-C310-jabita, TE-C310-sagam, TE-C310-kuli, TE-C310-udyogi

\chapter{Purse, List, Half, Daily wages, Employee}
\label{ch:te-purse-list-half-daily-wages-employee}
\hlchaptermodality{310}

\begin{chapteropening}
\textbf{By the end of this chapter:} \emph{I can say a purse, a list, half, daily wages and an employee.}

Complete the last lesson of chapter 310: I can say a purse, a list, half, daily wages and an employee.
\end{chapteropening}

\section[\texorpdfstring{\te{పర్సు} (parsu)}{పర్సు (parsu)}]{\te{పర్సు} (parsu) — a purse, a wallet}

\label{lesson:TE-C310-parsu}



\begin{quote}
Before the new one: say the Telugu for a water tap, then the Telugu for a flag.
\end{quote}

\subsection*{You'll want to know: \te{పర్సు}}
\textbf{\te{పర్సు}} — \emph{parsu} — \textquotedblleft{}a purse, a wallet\textquotedblright{}.

\begin{grammarlens}[title={What you've built}]
One more word for this chapter.
\end{grammarlens}

\subsection*{Your turn}
\begin{itemize}
\raggedright
  \item \emph{Say it:} \emph{parsu}
  \item \emph{Say it:} \emph{parsu}, once more
  \item \emph{Say it:} say \emph{parsu}
  \item \emph{Recall:} say the Telugu for a water tap, then the Telugu for a flag, then say \emph{parsu} again
\end{itemize}

\begin{tcolorbox}[breakable,colback=teal!4,colframe=teal!35!black,title={Before you move on}]
What does \te{పర్సు} mean? (\textquotedblleft{}A purse, a wallet\textquotedblright{}.) Say it once more.
\end{tcolorbox}
\section[\texorpdfstring{\te{జాబితా} (jābitā)}{జాబితా (jābitā)}]{\te{జాబితా} (jābitā) — a list}

\label{lesson:TE-C310-jabita}



\begin{quote}
Before the new one: say the Telugu for a flag, then the Telugu for a purse.
\end{quote}

\subsection*{You'll want to know: \te{జాబితా}}
\textbf{\te{జాబితా}} — \emph{jābitā} — \textquotedblleft{}a list\textquotedblright{}.

\begin{grammarlens}[title={What you've built}]
One more word for this chapter.
\end{grammarlens}

\subsection*{Your turn}
\begin{itemize}
\raggedright
  \item \emph{Say it:} \emph{jābitā}
  \item \emph{Say it:} \emph{jābitā}, once more
  \item \emph{Say it:} say \emph{jābitā}
  \item \emph{Recall:} say the Telugu for a flag, then the Telugu for a purse, then say \emph{jābitā} again
\end{itemize}

\begin{tcolorbox}[breakable,colback=teal!4,colframe=teal!35!black,title={Before you move on}]
What does \te{జాబితా} mean? (\textquotedblleft{}A list\textquotedblright{}.) Say it once more.
\end{tcolorbox}
\section[\texorpdfstring{\te{సగం} (sagaṁ)}{సగం (sagaṁ)}]{\te{సగం} (sagaṁ) — half}

\label{lesson:TE-C310-sagam}



\begin{quote}
Before the new one: say the Telugu for a purse, then the Telugu for a list.
\end{quote}

\subsection*{You'll want to know: \te{సగం}}
\textbf{\te{సగం}} — \emph{sagaṁ} — \textquotedblleft{}half\textquotedblright{}.

\begin{grammarlens}[title={What you've built}]
One more word for this chapter.
\end{grammarlens}

\subsection*{Your turn}
\begin{itemize}
\raggedright
  \item \emph{Say it:} \emph{sagaṁ}
  \item \emph{Say it:} \emph{sagaṁ}, once more
  \item \emph{Say it:} say \emph{sagaṁ}
  \item \emph{Recall:} say the Telugu for a purse, then the Telugu for a list, then say \emph{sagaṁ} again
\end{itemize}

\begin{tcolorbox}[breakable,colback=teal!4,colframe=teal!35!black,title={Before you move on}]
What does \te{సగం} mean? (\textquotedblleft{}Half\textquotedblright{}.) Say it once more.
\end{tcolorbox}
\section[\texorpdfstring{\te{కూలి} (kūli)}{కూలి (kūli)}]{\te{కూలి} (kūli) — daily wages; a labourer}

\label{lesson:TE-C310-kuli}



\begin{quote}
Before the new one: say the Telugu for a list, then the Telugu for half.
\end{quote}

\subsection*{You'll want to know: \te{కూలి}}
\textbf{\te{కూలి}} — \emph{kūli} — \textquotedblleft{}daily wages; a labourer\textquotedblright{}.

\begin{grammarlens}[title={What you've built}]
One more word for this chapter.
\end{grammarlens}

\subsection*{Your turn}
\begin{itemize}
\raggedright
  \item \emph{Say it:} \emph{kūli}
  \item \emph{Say it:} \emph{kūli}, once more
  \item \emph{Say it:} say \emph{kūli}
  \item \emph{Recall:} say the Telugu for a list, then the Telugu for half, then say \emph{kūli} again
\end{itemize}

\begin{tcolorbox}[breakable,colback=teal!4,colframe=teal!35!black,title={Before you move on}]
What does \te{కూలి} mean? (\textquotedblleft{}Daily wages; a labourer\textquotedblright{}.) Say it once more.
\end{tcolorbox}
\section[\texorpdfstring{\te{ఉద్యోగి} (udyōgi)}{ఉద్యోగి (udyōgi)}]{\te{ఉద్యోగి} (udyōgi) — an employee}

\label{lesson:TE-C310-udyogi}



\begin{quote}
Before the new one: say the Telugu for half, then the Telugu for daily wages.
\end{quote}

\subsection*{You'll want to know: \te{ఉద్యోగి}}
\textbf{\te{ఉద్యోగి}} — \emph{udyōgi} — \textquotedblleft{}an employee\textquotedblright{}.

\begin{grammarlens}[title={What you've built}]
One more word for this chapter.
\end{grammarlens}

\subsection*{Your turn}
\begin{itemize}
\raggedright
  \item \emph{Say it:} \emph{udyōgi}
  \item \emph{Say it:} \emph{udyōgi}, once more
  \item \emph{Say it:} say \emph{udyōgi}
  \item \emph{Recall:} say the Telugu for half, then the Telugu for daily wages, then say \emph{udyōgi} again
\end{itemize}

\begin{tcolorbox}[breakable,colback=teal!4,colframe=teal!35!black,title={Before you move on}]
What does \te{ఉద్యోగి} mean? (\textquotedblleft{}An employee\textquotedblright{}.) Say it once more.
\end{tcolorbox}
